\documentclass[11pt]{article}
\usepackage[utf8]{inputenc}
\usepackage[spanish]{babel}
\usepackage[margin=2.1cm]{geometry}
\usepackage{amsmath}
\usepackage{microtype}
\usepackage{xcolor}
\usepackage{booktabs}
\usepackage{array}
\usepackage[shortlabels]{enumitem}
\usepackage{titlesec}
\usepackage{hyperref}
\usepackage{fancyhdr}
\usepackage{parskip}

\definecolor{azulcf}{RGB}{25,55,110}
\definecolor{rojocf}{RGB}{170,35,45}
\definecolor{grisleve}{RGB}{90,90,90}

\hypersetup{colorlinks=true, linkcolor=azulcf, urlcolor=azulcf, citecolor=azulcf}

\titleformat{\section}{\large\bfseries\color{azulcf}}{\thesection.}{0.5em}{}
\titleformat{\subsection}{\normalsize\bfseries\color{rojocf}}{\thesubsection}{0.5em}{}
\titlespacing*{\section}{0pt}{1.1em}{0.5em}
\titlespacing*{\subsection}{0pt}{0.8em}{0.3em}

\setlength{\parskip}{0.35em}
\setlist{itemsep=1pt, topsep=2pt, parsep=0pt}

\pagestyle{fancy}
\fancyhf{}
\fancyhead[L]{\small\color{grisleve} CODEFEST AD ASTRA 2026 --- Etapa 1}
\fancyhead[R]{\small\color{grisleve} Base de Conocimiento}
\fancyfoot[C]{\small\thepage}
\renewcommand{\headrulewidth}{0.4pt}

\newcommand{\codigo}[1]{\texttt{\small #1}}
\newcommand{\seccionspec}[1]{\textit{(Sección #1 del spec)}}

\begin{document}

\begin{center}
{\LARGE\bfseries\color{azulcf} Documento Técnico --- Etapa 1}\\[0.3em]
{\large Construcción de la Base de Conocimiento}\\[0.2em]
{\normalsize\color{grisleve} CODEFEST AD ASTRA 2026 --- Equipo Laka}
\end{center}

\vspace{0.3em}
\hrule
\vspace{0.6em}

Este documento describe las decisiones de diseño tomadas para construir la base
de conocimiento vectorial exigida por la Etapa~1: extracción y limpieza de
fuentes, estrategia de chunking, encoder seleccionado, tipo de índice FAISS y
el módulo de recuperación que produce \codigo{resultados.jsonl}. Todas las
cifras citadas provienen de mediciones directas sobre el corpus, el índice y
la salida realmente entregados en \codigo{entrega/}, no de estimaciones.

\section{Arquitectura general}

El pipeline sigue el flujo de extremo a extremo de la Figura~1 del spec:
extracción y limpieza por fenómeno $\rightarrow$ chunking semántico
$\rightarrow$ codificación con encoder $\rightarrow$ indexación en FAISS
$\rightarrow$ recuperación. Las tres primeras etapas producen los artefactos
persistidos en el repositorio (\codigo{data/clean/*.jsonl},
\codigo{entrega/base\_vectorial/}); la última es \codigo{entrega/generador.py},
el script reproducible que consume esos artefactos y no construye nada nuevo.

\section{Extracción y limpieza}

Cada fenómeno tiene su propio script de extracción
(\codigo{src/f\{1,2,3\}\_limpieza\_extraccion.py}), autocontenidos a propósito
para no acoplar su evolución. Los tres comparten limpieza de caracteres de
control, reparación de guionado a fin de línea, eliminación de cabeceras/pies
repetidos y detección de idioma (\codigo{langdetect}).

Divergen donde el corpus lo exige: F1 y F3 extraen PDF con PyMuPDF en modo
\emph{texto plano} (una línea por línea visual, colapsando todo el
espaciado); F2 usa el modo \emph{blocks}, que preserva los saltos de párrafo
como separadores y añade OCR sobre imágenes sueltas (incluyendo un archivo
AVIF, con un reencodeo a PNG en memoria porque \codigo{pytesseract} no
decodifica AVIF directamente) y una reducción de catálogos de scraping a
\codigo{título | año | page\_url}. F3 añade lectura robusta de CSV con
detección de separador/codificación y extracción de teselas vectoriales
Mapbox (\codigo{.pbf}), leyendo capas y propiedades como pares
atributo:valor.

Un conjunto de documentos en idiomas fuera de español/inglés/portugués
(traducciones oficiales del mismo reporte publicadas por UNOOSA, SWF y SIPRI)
se excluye por una lista explícita de \codigo{doc\_id} verificados a mano, no
por \codigo{idioma\_detectado}: \codigo{langdetect} produce falsos positivos
frecuentes en texto corto o técnico (español marcado como catalán, alemán o
rumano en fuentes de CENIA y en datos \codigo{.pbf}) y en contenido no
lingüístico. Esos falsos positivos se retienen explícitamente en vez de
excluirse --- solo se excluyen los casos verificados de idioma genuinamente
distinto.

\section{Estrategia de chunking y justificación}
\seccionspec{3}

\subsection{Estrategia elegida: semántica con superposición}

Se implementó la estrategia \emph{semántica con superposición} (Sección~3.2):
en vez de trocear por tamaño fijo, se detectan cambios de tema calculando la
distancia coseno entre embeddings de oraciones consecutivas (cada oración se
embebe junto con una ventana de contexto de $\pm1$ oración, para estabilizar
el punto de corte) y se corta donde esa distancia supera el percentil~95 de
las distancias del propio documento --- un umbral adaptativo por documento, no
global. Los grupos semánticos resultantes por debajo de un mínimo de palabras
se fusionan con el siguiente para evitar fragmentación excesiva.

\subsection{Restricciones duras}

Sobre el agrupamiento semántico se imponen dos restricciones que nunca se
relajan:

\begin{itemize}
\item \textbf{Completitud lingüística} \seccionspec{3.3}: los cortes solo
ocurren en fronteras de oración (tokenización con NLTK, adaptada por idioma
detectado). Para formatos tabulares (csv/xlsx/pbf) se usa el salto de línea
como frontera en vez de NLTK: Punkt trata un punto tras un número como parte
del número, así que un CSV completo se tokenizaba como una sola oración
gigante.
\item \textbf{Límite de 250 palabras} \seccionspec{9.2.1}: un grupo semántico
que excede el límite se subdivide empaquetando oraciones, igual que un
chunking por tamaño fijo, pero solo como mecanismo de contención --- nunca
como estrategia primaria.
\end{itemize}

\subsection{Caso patológico: oraciones que exceden el límite por sí solas}

Documentos de política pública con cláusulas largas, o texto tabular mal
segmentado por el tokenizador, producen ``oraciones'' que superan las 250
palabras sin ningún punto de corte oracional disponible. Una heurística de
densidad de conectores gramaticales (es/pt/en) distingue prosa genuina de
artefactos no lingüísticos: la prosa se corta en el separador de cláusula
(coma, punto y coma o dos puntos) más cercano al límite; los artefactos se
cortan por bloques de palabras. Cada partición forzada se registra en un
archivo de auditoría para trazabilidad.

Como red de seguridad final, cada chunk se re-verifica contra el
\emph{tokenizer real} del encoder (no un proxy por palabras): el español
técnico infla hasta 2$\times$ en subword tokens, así que un chunk de
$\le$250 palabras puede en teoría exceder la ventana de 512 tokens del
encoder. Si ocurre, se vuelve a partir por conteo de tokens exacto antes de
escribirse en la metadata.

\subsection{Resultado medido}

\begin{center}
\begin{tabular}{@{}lr@{}}
\toprule
\textbf{Métrica} & \textbf{Valor} \\
\midrule
Chunks totales & 330{,}416 \\
Documentos únicos (\codigo{doc\_id}) & 1{,}780 \\
Palabras por chunk --- mínimo / mediana / p90 / máximo & 1 / 140 / 228 / 250 \\
Chunks que exceden 250 palabras & \textbf{0} \\
\bottomrule
\end{tabular}
\end{center}

Ningún chunk excede el límite: la restricción dura de la Sección~9.2.1 se
cumple por construcción, no por post-filtrado. La distribución por fenómeno
es F1~=~200{,}780, F2~=~51{,}859, F3~=~77{,}777 chunks, sobre 1{,}780
documentos únicos (454 en F1, 449 en F2, 877 en F3); la distribución por
formato muestra que más de la mitad del corpus (53\,\%) proviene de formatos
tabulares (csv, xlsx, pbf).

\section{Encoder y criterios de elección}
\seccionspec{4}

\subsection{Encoder seleccionado}

\codigo{intfloat/multilingual-e5-small} (familia BERT, 117.65\,M de
parámetros, licencia MIT, dimensión 384, ventana de 512 tokens). Requiere
prefijar el texto con
\codigo{"query: "} (consultas) o \codigo{"passage: "} (documentos) antes de
codificar --- aplicado consistentemente en indexación y en recuperación, y
verificado explícitamente (Sección~\ref{sec:faiss}).

\subsection{Criterios (Sección 4.3)}

\begin{itemize}
\item \textbf{Soporte multilingüe nativo}: el corpus mezcla español, inglés y
portugués. Se verificó con una prueba de humo de tres consultas (una por
fenómeno, cada una en un idioma distinto) que recupera correctamente entre
idiomas: 5/5 resultados del fenómeno esperado en los tres casos.
\item \textbf{Dimensionalidad moderada} (384): equilibrio entre costo de
almacenamiento/búsqueda y expresividad, adecuado para un índice
\codigo{IndexFlatIP} de 330 mil vectores servido en CPU.
\item \textbf{Longitud máxima de entrada} (512 tokens): compatible con el
límite de 250 palabras de los chunks, con margen para la inflación de
subword tokenization verificada en el paso de chunking.
\item \textbf{Licencia} MIT: sin restricciones para el uso en el reto.
\item \textbf{Eficiencia computacional}: modelo ``small'' de la familia E5,
dimensionado para correr en CPU dentro de los tiempos del reto (índice
completo codificado sin GPU).
\end{itemize}

\subsection{Comparación empírica}

\codigo{src/benchmark\_modelos.py} implementa una comparación empírica de
cinco encoders candidatos contra el ground truth anotado del equipo,
calculando NDCG@10 y F1@3 y combinándolos por Conteo de Borda
(Sección~11.2):

\begin{itemize}[itemsep=0pt]
\small
\item \codigo{paraphrase-multilingual-MiniLM-L12-v2}
\item \codigo{multilingual-e5-small} (seleccionado)
\item \codigo{paraphrase-multilingual-mpnet-base-v2}
\item \codigo{BAAI/bge-m3}
\item \codigo{multilingual-e5-base}
\end{itemize}

\textbf{Nota de honestidad}: la corrida completa sobre los cinco modelos no
se completó ni quedó persistida en este repositorio --- solo sobrevive el
archivo de resultados de un modelo (\codigo{paraphrase-multilingual-MiniLM
-L12-v2}), sin tabla comparativa final. La elección de
\codigo{multilingual-e5-small} se sostiene, por tanto, en los criterios
cualitativos de la Sección~4.3 más la verificación multilingüe empírica de
la Sección~\ref{sec:faiss}, no en un ranking numérico de Borda entre los
cinco candidatos.

No se usaron múltiples encoders combinados (Sección~4.4): un único encoder
multilingüe ya cubre los tres idiomas del corpus sin necesitar fusión de
rankings (CombSUM/CombMNZ/RRF), lo que simplifica el pipeline de recuperación
sin sacrificar cobertura de idioma.

\section{Índice FAISS}
\label{sec:faiss}
\seccionspec{5}

\codigo{IndexFlatIP} sobre vectores L2-normalizados (equivalente a similitud
coseno exacta, Sección~5.2), con 330{,}416 vectores de dimensión 384. Es la
opción recomendada por el spec para el volumen de este corpus: búsqueda
exacta sin pérdida de precisión, sin la complejidad de ajustar
\codigo{nprobe} (IVF) o el uso de memoria de un grafo HNSW. El archivo
resultante pesa $\sim$508\,MB; la codificación completa del corpus tomó
9.5 minutos sobre GPU ($\sim$580 chunks/segundo). El orden de inserción se
mantuvo idéntico al orden de lectura de \codigo{metadata.jsonl}, garantizando
que el id interno de FAISS coincida con el número de línea del almacén de
metadata (Sección~1.4).

Verificaciones realizadas sobre el índice entregado
(\codigo{src/validar\_indice.py}):

\begin{center}
\begin{tabular}{@{}p{5.3cm}p{9.8cm}@{}}
\toprule
\textbf{Chequeo} & \textbf{Resultado} \\
\midrule
Alineación metadata $\leftrightarrow$ índice & 330{,}416 líneas == 330{,}416 vectores; la línea $N$ corresponde al id interno $N$ \\
Normalización & Muestra de 200 vectores, norma unitaria en todos \\
Identidad vector$\leftrightarrow$texto & 30 muestras: similitud coseno mínima $=1.000000$ entre el vector guardado y el reembebido desde el texto \\
Humo multilingüe (es/en/pt) & 5/5 resultados del fenómeno esperado en cada una de las tres consultas \\
\bottomrule
\end{tabular}
\end{center}

El chequeo de identidad es la prueba más exigente: reconstruye el vector en
una posición y lo compara contra el reembebido desde el texto de esa misma
línea de metadata. Una similitud de 1.000000 descarta cualquier
desalineación entre el orden de indexación y el orden del archivo de
metadata.

\section{Módulo de recuperación}
\seccionspec{8}

\codigo{entrega/generador.py} implementa el flujo de la Sección~8.1 sin usar
en ningún punto un modelo generativo (Sección~8.3): la consulta se codifica
con el mismo encoder e idéntico prefijo E5, se normaliza, y se busca en
FAISS. Toda decisión opera sobre vectores, puntuaciones coseno y metadata.

\subsection{Lectura de consultas}

El script acepta el archivo de consultas en varios formatos
(\codigo{.jsonl}/\codigo{.json}/\codigo{.csv}/\codigo{.xlsx}), autodetectado
por extensión y con alias flexibles de nombre de columna
(\codigo{query\_id}/\codigo{id}/\codigo{qid}, \codigo{query}/\codigo{pregunta}/
\codigo{texto}). El archivo \codigo{entrega/\allowbreak consultas.jsonl} se
generó a partir del PDF oficial de los organizadores
(\codigo{Extracto\_\allowbreak Preguntas\_50\_v2.pdf}, en \codigo{data/raw/})
con un parser por expresión regular sobre el patrón \codigo{qNNN}, verificado
línea por línea contra la lista de referencia usada durante el desarrollo:
50/50 idénticas.
Se generó como paso explícito de preparación --- necesario para que
\codigo{generador.py} sea ejecutable de punta a punta, ya que el spec exige
en la Sección~1.4 que el script ``lea el archivo de consultas'' sin fijar su
formato, y sin él la Sección~1.4 también advierte que una entrega no
reproducible se excluye de la evaluación.

\subsection{Búsqueda y agregación a documento (Sección 8.6)}

El tamaño del pool de candidatos $k$ es dinámico: arranca en 30 y se duplica
hasta 200 si hace falta, garantizando que toda consulta obtenga exactamente
3 documentos y 10 fragmentos --- requisito estricto de la Sección~9.3.2,
donde un conteo incorrecto se penaliza o descarta. Los candidatos se agrupan
por \codigo{doc\_id} y se agregan con \textbf{max pooling} (el score del
fragmento más relevante representa al documento), una de las tres
estrategias que nombra textualmente la Sección~8.6.

Esta elección no fue arbitraria: durante la validación se detectó un sesgo
sistemático real, documentado en la Sección~\ref{sec:agregacion}, y se
probaron y descartaron dos alternativas antes de confirmar max pooling puro
como la configuración final.

\subsection{Límite de 250 palabras por fragmento (Sección 9.2.1)}

Los chunks recuperados que superan 250 palabras se dividen en sub-fragmentos
mediante tokenización de oraciones (NLTK) empaquetadas de forma voraz hasta
el límite, preservando siempre fronteras oracionales completas; todas las
sub-piezas de un mismo chunk comparten el mismo \codigo{chunk\_id} y ocupan
posiciones de ranking independientes, como exige el spec. En la práctica
esta ruta \textbf{no se activa} sobre el corpus actual (ningún chunk supera
250 palabras, Sección~3), pero se implementó y validó con texto sintético
para no dejar un caso límite sin cubrir.

\subsection{Autovalidación de salida}

\codigo{generador.py} valida su propia salida al terminar: cuenta líneas de
\codigo{resultados.jsonl}, exige exactamente 3 documentos y 10 fragmentos por
consulta, y verifica que ningún fragmento exceda el límite de palabras ---
una red de seguridad barata contra errores de esquema antes de la entrega.

\subsection{Observaciones verificadas sobre la salida entregada}
\label{sec:observaciones}

Dos comportamientos medidos directamente sobre \codigo{resultados.jsonl},
relevantes para interpretar la Sección~\ref{sec:evaluacion}:

\begin{itemize}
\item \textbf{Sin normalización de texto}: el 89\,\% de los fragmentos
(445/500) arrastra saltos de línea de la extracción de PDF a mitad de
frase, mientras que el 100\,\% de los excertos anotados en el ground truth
de referencia son de una sola línea. Como la relevancia se juzga por el
\emph{contenido} del campo \codigo{text} (Sección~10.2.1), esto es una
oportunidad de mejora de bajo costo: normalizar espacios no infringe el
esquema (la Sección~9.2.1 ya autoriza modificar el texto de salida) y no
toca el campo \codigo{texto} de \codigo{metadata.jsonl}, que la Tabla~1
exige sin modificaciones.
\item \textbf{Sin deduplicación}: al no filtrar por similitud, 19 de las 50
consultas (38\,\%) tienen al menos dos fragmentos casi idénticos dentro de
su propio top-10 --- esperable dado que el corpus incluye documentos con
filas de contenido muy repetitivo. La Sección~8.7 autoriza deduplicar por
similitud entre vectores o por campos de metadata; no se implementó en esta
entrega por restricción de tiempo.
\end{itemize}

\section{Evaluación y limitaciones}
\label{sec:evaluacion}
\seccionspec{10}

\subsection{Experimentos de agregación: hipótesis probadas y descartadas}
\label{sec:agregacion}

Durante la validación se detectó un sesgo sistemático: documentos con
muchísimos chunks (p.\,ej. el Índice Latinoamericano de IA, hasta 975 chunks
frente a una mediana de 7 en todo el corpus) dominaban el top-3 de consultas
específicas incluso cuando existía un documento más pequeño y preciso con el
contenido correcto --- confirmado con la consulta \emph{``¿Qué desafíos
plantea el uso de inteligencia artificial en las operaciones espaciales?''},
donde el sistema devolvió 0/3 documentos del fenómeno correcto pese a que
existían 86 chunks de contenido de IA aplicada al dominio espacial en el
corpus. La causa raíz: con max pooling, el score de un documento es el
máximo entre sus chunks candidatos --- el mejor resultado de $N$ intentos---,
y un documento con más chunks tiene estadísticamente más oportunidades de
que uno de ellos puntúe alto por azar (estadística de valores extremos), sin
que eso implique mayor relevancia real.

Se probaron tres alternativas contra las 50 consultas reales (no contra
casos aislados), usando como proxy la coincidencia de fenómeno entre
documento y consulta, validado además con casos de depuración individuales
antes de la corrida completa:

\begin{center}
\small
\begin{tabular}{@{}p{2.3cm}p{3.1cm}p{6.0cm}p{1.4cm}@{}}
\toprule
\textbf{Hipótesis} & \textbf{Mecanismo} & \textbf{Resultado empírico (50 consultas)} & \textbf{Decisión} \\
\midrule
Penalización log & $\text{score} - \beta\log(n)$ & F1: 85\%$\to$48\% (penaliza injustamente a
documentos de tamaño típico de F1 frente a F3, mediana 2 vs.\ 14 chunks/doc) & Descartada \\
+ umbral (100) & $\text{score} - \beta\max(0,\log n - \log 100)$ & Sin cambio real (F1 igual en 48\%); total agregado 89.3\%$\to$78.7\% (\,$-$10.6\,pp) & Descartada \\
Media de top-$k$ & promedio de los 3 chunks más altos & Empeoró el caso diagnóstico: premió amplitud
genérica sobre enfoque específico y certero & Descartada \\
\bottomrule
\end{tabular}
\end{center}

Las tres empeoraron el resultado agregado frente a max pooling sin
modificar, así que se revirtieron. La configuración final entregada usa max
pooling puro.

\subsection{Validación contra ground truth}

El archivo \codigo{FASE ORDENADA CODEFEST.xlsx}, provisto por el equipo
organizador, contiene relevancia verificada (pregunta, fragmento y documento
fuente) para \textbf{7 de las 50 consultas} (q002, q003, q004, q033--q036;
hojas F1 y F3, sin hoja para F2), cruzado contra el \codigo{doc\_id} interno
vía coincidencia normalizada del nombre de archivo (campo \codigo{fuente}).
Es una muestra pequeña que no reemplaza el NDCG@10/F1@3 oficial sobre las 50
consultas --- cuyo ground truth no es público durante el reto --- pero es
evidencia real medida dos veces, con dos metodologías independientes:

\begin{itemize}
\item \textbf{Proxy de coincidencia} (top-3/top-10, sin ponderar por
posición): 1/7 consultas (14\,\%) obtuvieron el documento correcto en el
top-3; 3/7 (43\,\%) lo obtuvieron en el top-3 o entre los 10 fragmentos.
\item \textbf{NDCG@10/F1@3 exactos} (Sección~10.2, recalculados con
emparejamiento uno-a-uno entre fragmento y excerto para evitar que el DCG
supere al IDCG): $\overline{\text{NDCG@10}} = 0.1429$,
$\overline{\text{F1@3}} = 0.0714$ sobre las mismas 7 consultas.
\end{itemize}

Ambas metodologías apuntan en la misma dirección --- desempeño débil sobre
esta muestra pequeña, con margen real de mejora --- lo cual da más confianza
que cualquiera de las dos por separado. El patrón de falla de la
Sección~\ref{sec:agregacion} se confirma en esta muestra independiente: las
consultas que combinan el término genérico ``inteligencia artificial'' con
un calificador técnico específico (``sistemas no tripulados'', ``conflictos
recientes'', ``operaciones espaciales'') tienden a recuperar el Índice
Latinoamericano de IA --- un reporte extenso y genérico --- en lugar de la
fuente específica correcta (informes DAIO, CSIS), incluso cuando esta última
existe en el corpus con contenido directamente relevante. Se identifica como
\textbf{limitación conocida} del enfoque de recuperación puramente vectorial
con un encoder compacto (117\,M parámetros): el término genérico domina la
señal semántica sobre el calificador específico.

Con una muestra de 7 consultas cualquier cifra puntual es sensible a 1--2
casos individuales; estos números deben leerse como evidencia direccional,
no como una estimación precisa de NDCG@10/F1@3 sobre las 50 consultas
oficiales.

\subsection{Trabajo futuro propuesto}

La dirección con mayor evidencia de impacto es \textbf{recuperación
híbrida} (búsqueda léxica tipo BM25 combinada con la búsqueda vectorial
actual vía Reciprocal Rank Fusion, Sección~8.4): un match léxico exacto de
términos técnicos específicos no depende del tamaño del documento ni de
matices que un encoder compacto no capture, y ataca directamente la causa
raíz identificada en la Sección~\ref{sec:agregacion}. No se implementó en
esta entrega por restricción de tiempo; queda documentado como el próximo
paso concreto. Las dos observaciones de la Sección~\ref{sec:observaciones}
(normalización de texto de salida y deduplicación) son mejoras de menor
esfuerzo, con margen explícito en la Sección~8.7 del spec, que podrían
implementarse antes.

\section{Grafo de conocimiento}
\seccionspec{7}

No se implementó el componente bonus de grafo de conocimiento en esta
entrega.

\end{document}
